\section{Joint policy-cost and resource decisions}\label{sec:joint50}
Let a finite catalogue $\mathcal P$ contain implemented policies, including incumbents, repaired policies and conventional alternatives. Construction, verification and direct evaluation are distinct charged activities. Denote their measured deployment-relevant resource vector by $r_j$; the published timing convention must specify whether an evaluation shared by several candidates is a common fixed charge or an attributable marginal charge. In particular, a joint paired evaluation cannot be counted as zero or silently charged only to the preferred generator.

Conditional on construction information, suppose a simultaneous event satisfies
$J_j\in[L_j,U_j]$ for every $j\in\mathcal P$. For nonnegative prices $\theta$, let $k_j(\theta)$ be a known implementation charge, for example $\theta\cdot r_j$ plus an information fee. Define
\begin{equation}\label{eq:jointselector50}
 \widehat j(\theta)\in\arg\min_{j\in\mathcal P}
   \{U_j+k_j(\theta)\},\qquad
 R(\theta)=U_{\widehat j}+k_{\widehat j}(\theta)
            -\min_{j\in\mathcal P}\{L_j+k_j(\theta)\}.
\end{equation}
A fixed lexicographic tie rule makes the selection measurable.
\begin{proposition}[Actual net-cost frontier]\label{prop:joint50}
On the simultaneous event, for every nonnegative price vector, including one selected after observing the intervals,
\[
 0\le J_{\widehat j}+k_{\widehat j}(\theta)
       -\min_j\{J_j+k_j(\theta)\}\le R(\theta).
\]
A proposed replacement of incumbent $i$ by $j$ is certified nonworsening in net cost whenever a simultaneous upper bound on $J_j-J_i$, plus $k_j(\theta)-k_i(\theta)$ and the replacement fee, is nonpositive. Otherwise a rule retaining $i$ remains safe. Policy-specific certificate widths can be displayed alongside this frontier, but do not determine its ordinate.
\end{proposition}
\begin{proof}
Bound the selected true cost above by its upper line and the best catalogue cost below by the minimum lower line. All inequalities hold on the same finite event for every price, so a new price query does not consume another probability allocation. The replacement assertion applies the paired upper bound directly.\end{proof}

Absolute intervals can be considerably wider than centered paired differences. For a decision between an incumbent and its final repair, the centered interval is therefore the appropriate replacement input. A statistically unresolved interval is not evidence of equivalence. In contrast, Corollary~\ref{cor:terminal50} supplies a deterministic nonincrease, although it does not necessarily certify a strictly positive improvement large enough to cover an implementation fee.

\subsection{Information technology and the price of conservatism}\label{sec:technology50}
An explicit theoretical observation technology gives the bit price units. Suppose a device reports a dyadic cell in each coordinate at date $t$. With $b$ bits per coordinate it transmits $db$ bits and requires $\tau db$ resource units, where $\tau$ is a technological conversion factor. At unit resource price $w$, and with a setup cost $K_b$, its present cost is
\begin{equation}\label{eq:technology50}
 C_b(w,\tau)=K_b+w\tau db\sum_{t<T}B_t.
\end{equation}
Each reported cell entails its own feasible action contract. In the retained sensor experiment, $K_b=0$ and the normalized price $\lambda=w\tau$ takes the declared values. This supplies a structural interpretation and a sensitivity variable, not an estimate of actual device costs. The dataset contains no empirical measurements from which $w$ or $\tau$ could be calibrated.

Theorem~\ref{thm:net49} and Proposition~\ref{prop:joint50} apply to the finite bit catalogue with charges \eqref{eq:technology50}. Their upper-minus-lower envelope is a valid bound on the cost of conservative choice relative to the unknown true catalogue minimizer. Replacing that minimizer by the lowest observed sample mean and calling it the ex post true optimum is not justified. The retained exact price envelopes and their regret bounds answer the identified decision question without making this substitution. A richer technology with fixed charges, nonlinear bandwidth costs or precision-dependent delays can be assessed by changing the known charge functions, but new unobserved policies require new cost bounds.

A similarly explicit interpretation applies to terminal repair. A controller that incurs cost $K$ to install the additional gate, and $q$ per deployed final action, is worth adopting only when its expected discounted policy gain covers those charges. The published centered gain interval defines the entire break-even range. Its mathematical safety does not by itself imply economic adoption at positive computation prices.
